\section*{Chapitre \ref{ch:s2:skew-and-term-structure-relative-value} --- Relative value sur skew et structure par terme}

\begin{solution}{exo:s2:skew-and-term-structure-relative-value:1}
$100 - 10 \times (-2.3) = 123$.
\end{solution}

\begin{solution}{exo:s2:skew-and-term-structure-relative-value:2}
La volatilité baisse de $0.4/20 = 2~\%$ de sa propre valeur, donc la position courte rapporte 2 unités.
\end{solution}

\begin{solution}{exo:s2:skew-and-term-structure-relative-value:3}
$\ln 0.5/\ln 0.98 = 34.3$ jours de bourse.
\end{solution}

\begin{solution}{exo:s2:skew-and-term-structure-relative-value:4}
La volatilité revient à la moyenne, si bien que la structure par terme juste s'inverse lorsque la volatilité est élevée et présente une pente ascendante lorsqu'elle est faible~: la pente brute est corrélée au niveau à
$-0.83$ sur la surface synthétique et $-0.54$ sur la courbe réelle du VIX. Acheter le calendar spread lorsque la pente est faible revient à l'acheter lorsque la volatilité est élevée, un pari
que la volatilité baisse. Le résidu sur le niveau élimine l'essentiel de cet effet.
\end{solution}

\begin{solution}{exo:s2:skew-and-term-structure-relative-value:5}
Comme signal de timing pour les trades de variance courte du chapitre~1~: vendre de la variance lorsque la pente indique que la prime de variance est élevée. Pour un trader de calendar spread
ce résultat montre que la pente traduit surtout le prix du risque de variance, et non une prévision de volatilité, si bien qu'un calendar spread signalé par la pente est exposé à la prime
de variance, soit le trade de niveau qu'il devait éviter.
\end{solution}

\begin{solution}{exo:s2:skew-and-term-structure-relative-value:6}
Vendre le skew de put de l'indice et acheter les skews de put des membres (un risk reversal sur l'indice contre des risk reversals sur les membres), neutre en vega et
couvert en delta. Son principal risque est la corrélation~: lors d'un krach, chaque membre baisse avec l'indice et la cherté du skew d'indice se trouve justifiée, ce qui est
le risque de dispersion du chapitre~2 sous forme de skew.
\end{solution}

\begin{solution}{exo:s2:skew-and-term-structure-relative-value:7}
Le book de skew brut affiche encore un ratio de Sharpe de 1.08 (pire mois $-2.8$ écarts-types)~: il n'a jamais eu besoin de l'effet injecté, car il
vend l'aile des puts lorsque la volatilité est élevée. Le book de skew résiduel perd ($-0.41$, un pire mois de $-11.7$ écarts-types), puisque son signal est désormais
qu'à l'erreur de la régression. Les deux calendar spreads perdent ($-0.12$ brut, $-0.11$ résiduel) à cause des coûts et du gamma.
\end{solution}

\begin{solution}{exo:s2:skew-and-term-structure-relative-value:8}
L'attribution dit autre chose. Seulement 9~\% du P\&L de ce book provient du retour du skew, 94~\% du spot et du temps~; le skew brut est pentu lorsque la volatilité est
élevée, si bien que le book vend l'aile des puts lorsque la volatilité est élevée. Avec le bruit injecté désactivé, il rapporte encore 1.08. Le ratio de Sharpe est réel, mais il
rémunère la prime de variance et le risque de krach qui l'accompagne, et non le retour à la moyenne du skew.
\end{solution}

\begin{solution}{pb:s2:skew-and-term-structure-relative-value:1}
\begin{enumerate}
\item Un \omterm{def:s2:skew-and-term-structure-relative-value:skew}{trade de skew}~: positions longues et courtes sur options à deux prix d'exercice d'une même échéance, neutres en vega et couvertes en delta, sur l'écart entre leurs volatilités. Un
\omterm{def:s2:skew-and-term-structure-relative-value:term}{trade de structure par terme}~: le même sur deux échéances.
\item $100 - 10S$, où $S$ désigne le prix de l'asymétrie du S\&P 500~; 123.2 en moyenne depuis 1990, 110.0 à 147.3 pour les 90~\% centraux, 101.23 au plus bas
(19 mars 1991) et 183.12 au plus haut (18 février 2025).
\item Les distributions risque-neutre des actions individuelles sont nettement moins asymétriques négativement que celle de l'indice.
\item Le skew juste vaut $-0.12$ au niveau moyen et varie de $-0.4$ par unité de volatilité à un mois~; les bruits sont un AR(1) sur le skew (écart-type 0.06) et sur
la volatilité à trois mois (1.5 points de volatilité), avec une demi-vie de 34 jours.
\item 1.93 points en moyenne~; inversée 7.6~\% des jours~; $-18.23$ le 12 mars 2020.
\item La volatilité revient à la moyenne, si bien qu'une volatilité élevée inverse la courbe~; la corrélation vaut $-0.54$.
\item La pente porte le prix du risque de variance et prédit les rendements des swaps de variance, des contrats à terme sur le VIX et des straddles.
\item RVX 5.2 points au-dessus du VIX (écart-type 2.3, demi-vie 15.7 jours)~; VXN 3.0 au-dessus (écart-type 2.3, demi-vie 15.3).
\item Niveau, pente et skew extraits de la surface~; pente et skew régressés sur le niveau avec fenêtres croissantes, en conservant les résidus.
\item Une unité par hausse de 1~\% de la volatilité de la jambe à l'entrée en position~; P\&L exprimé en pourcentage de volatilité.
\item Réévaluation dans l'ordre~: spot et temps, puis niveau, puis bruit du skew, puis bruit de la pente~; les compartiments s'additionnent, et déplacer le niveau en premier
range les variations justes du skew et de la pente dans le compartiment du niveau.
\item 1.13, 0.75, 0.07 et 0.64~; spot et temps 94, 73, $-31$ et 97~\%~; niveau 5, 24, 160 et $-16$~\%~; skew ou pente 9, 16, 155 et 51~\%.
\item \textbf{Résultat nommé.} Les \omterm{def:s2:skew-and-term-structure-relative-value:skew}{trades de skew} et de calendar spread signalés sur les résidus rapportent des ratios de Sharpe de 0.75 et 0.64 sur cette graine (0.39 et 0.38 en
moyenne sur huit). Seulement 16~\% et 51~\% de leur P\&L proviennent du skew et de la pente injectés~; le book de skew brut rapporte 1.13 avec 9~\%, et le
calendar spread brut a un P\&L de 160~\% de niveau.
\item Ratios de Sharpe de 1.01 (compartiment du skew) et 1.51 (compartiment de la pente).
\item Moyennes de 0.90, 0.39, $-0.07$ et 0.38~; cette graine est du côté favorable pour les books résiduels.
\item Le risk reversal~: 12.0 unités par mois venant du spot et du temps, $-15.9$ venant du niveau, $-6.2$ au total. Le calendar spread~: 9.6, $-3.2$, 4.7.
\item Position longue de deux unités d'aile de calls avant une hausse de 4.5~\% avec la volatilité à un mois passant de 15.2~\% à 9.7~\%~: $-289$ unités, 4.5 écarts-types.
\item Le book de skew brut rapporte encore 1.08~; les books résiduels et les calendar spreads perdent.
\item Attribuer le P\&L, comparer les signaux bruts et résiduels, inclure les mouvements conjoints du spot et de la volatilité, imputer le coût de l'option, rapporter plusieurs graines.
\item Tout ce que ses jambes ne compensent pas~: les mouvements conjoints du spot et de la volatilité, et le gamma du mois proche.
\end{enumerate}
\end{solution}

\begin{solution}{iq:s2:skew-and-term-structure-relative-value:1}
Position longue sur un call hors de la monnaie et courte sur un put hors de la monnaie (ou l'inverse), à deltas ou vegas égaux. Son prix, exprimé comme la volatilité du call moins
celle du put, mesure le skew~: une valeur négative marquée indique que l'aile des puts est chère.
\end{solution}

\begin{solution}{iq:s2:skew-and-term-structure-relative-value:2}
La volatilité monte et la courbe s'inverse~: la volatilité de la jambe courte à un mois monte davantage que celle de la jambe longue à trois mois, si bien que le mouvement de niveau perd~; la jambe courte proche a
plus de gamma que la jambe longue éloignée, si bien que les grands mouvements perdent aussi. Les deux pénalisent.
\end{solution}

\begin{solution}{iq:s2:skew-and-term-structure-relative-value:3}
Régresser le skew sur le niveau et vérifier la corrélation~; trader le résidu autant que le signal brut et comparer~; attribuer le P\&L entre les compartiments de niveau et
de skew. Si le P\&L du signal brut vient surtout du spot, du temps ou du niveau, c'est un signal de niveau.
\end{solution}

\begin{solution}{iq:s2:skew-and-term-structure-relative-value:4}
Gamma et theta qui ne se compensent pas~; vega qui apparaît avec les mouvements du spot (une aile courte qui passe à la monnaie)~; le mouvement conjoint du spot, de la volatilité et du skew~; sauts
entre les couvertures~; la base entre sous-jacents~; les coûts de roll et de couverture.
\end{solution}

\begin{solution}{iq:s2:skew-and-term-structure-relative-value:5}
Réévaluer chaque position selon un ordre fixe des mouvements de marché (spot et temps, puis niveau, puis paramètres de forme) et enregistrer la variation de chaque étape~; le dernier
état doit égaler la surface complète du lendemain, si bien que les compartiments s'additionnent au total par construction. La surface doit être une fonction d'un petit nombre de
paramètres, assortie d'une règle sur sa variation avec le spot.
\end{solution}

\begin{solution}{iq:s2:skew-and-term-structure-relative-value:6}
Sur $\mathrm{d}t$, la valeur de l'option couverte varie de $\Theta\,\mathrm{d}t + \tfrac12\Gamma(\mathrm{d}S)^2$, et Black--Scholes donne $\Theta =
-\tfrac12\Gamma S^2\sigma_i^2$~; avec $(\mathrm{d}S)^2 \approx S^2\sigma_r^2\,\mathrm{d}t$, le P\&L vaut $\tfrac12\Gamma S^2(\sigma_r^2 -
\sigma_i^2)\,\mathrm{d}t$. Le vega vaut $S\varphi(d_1)\sqrt\tau$ et le gamma $\varphi(d_1)/(S\sigma_i\sqrt\tau)$, donc le gamma par unité de vega est
$1/(S^2\sigma_i\tau)$~: la jambe à volatilité implicite plus élevée a moins de gamma par unité de vega.
\end{solution}
